\documentclass{article}
\usepackage{/globals/laws}
\usepackage{/globals/de}
\usepackage{/globals/anki}
\ankiprefix{staaten}
\ankitopic{Staaten}
\ankishow
\begin{document}
\section{Staaten}
Ein Staat liegt nach Jellinek vor wenn drei Vorraussetzungen erfüllt sind; es liegt ein Staatsgebiet, ein Staatsvolk und eine Staatsgewalt vor.
\begin{description}
 \item[Staatsgebiet] Ein bestimmter, begrenzter Teil der Erdoberfläche, auf dem längerzeitig Menschen leben können.
 \item[Staatsvolk] Staatsangehörige Personen die dauerhaft verbunden sind.
 \item[Staatsgewalt] Tatsächliche Herrschaftsmacht im Staatsgebiet über das Staatsvolk, die nicht durch eine höhere Autorität abgeleitet wurde.
\end{description} 

\ankinote{merkmale}{Merkmale}{Staatsgebiet, Staatsvolk, Staatsgewalt}

\subsection{Staatsrecht}
Das Staatsrecht ist teil des öffentlichen Rechts und beinhaltet das Verfassungsrecht. Es behandelt alles rechtliche zum Verhältnis des Bürgers zum Staat, zu der Organisation und den Prinzipien des Staates, so wie diese im Grundgesetz und anderen Gesetzes geregelt ist.

\ankinote{staatsrecht-pos}{Staatsrecht Position}{Zwischen Öffentliches Recht und Verfassungsrecht}
\ankinote{staatsrecht-inhalt}{Staatsrecht Inhalt}{Organisation und Prinzipien des Staates nach Grundgeset etc}
\end{document}